\documentclass[11pt,a4paper]{article}

% ── Codificación, fuentes y lenguaje ───────────────────────────────────────────
\usepackage[utf8]{inputenc}
\usepackage[T1]{fontenc}
\usepackage[default,scale=0.95]{sourcesanspro}
\renewcommand{\familydefault}{\sfdefault}
\usepackage[spanish,es-noshorthands,es-tabla]{babel}

% ── Geometría y espaciado ─────────────────────────────────────────────────────
\usepackage[top=2.2cm, bottom=2.5cm, left=2.2cm, right=2.2cm, headheight=15pt]{geometry}
\usepackage{setspace}
\setstretch{1.18}
\usepackage{parskip}

% ── Matemáticas y unidades ────────────────────────────────────────────────────
\usepackage{amsmath}
\usepackage{amssymb}
\usepackage{siunitx}

% ── Circuitos y gráficos ──────────────────────────────────────────────────────
\usepackage[european, straightvoltages]{circuitikz}
\usepackage{tikz}
\usetikzlibrary{arrows.meta, positioning, shapes.geometric, fit, backgrounds}

% ── Tablas ────────────────────────────────────────────────────────────────────
\usepackage{booktabs}
\usepackage{tabularx}
\usepackage{array}
\usepackage{longtable}
\usepackage{colortbl}
\usepackage{enumitem}
\usepackage{url}

% ── Colores, cajas e iconos ───────────────────────────────────────────────────
\usepackage[dvipsnames]{xcolor}
\usepackage{tcolorbox}
\tcbuselibrary{skins, breakable}
\usepackage{fontawesome5}
\usepackage{graphicx}

% ── Listados de código (dark-mode infográfico) ────────────────────────────────
\usepackage{listings}

% ── Secciones con barra de color (infografía) ────────────────────────────────
\usepackage{titlesec}
\titleformat{\section}
  {\normalfont\LARGE\bfseries\color{azulNoche}}
  {\colorbox{cyanNeon}{\color{fondoOscuro}\thesection}}{0.7em}{}
  [\vspace{2pt}\color{cyanNeon}\rule{\linewidth}{1.8pt}]
\titlespacing*{\section}{0pt}{22pt}{9pt}

\titleformat{\subsection}
  {\normalfont\large\bfseries\color{azulNoche}}
  {\textcolor{cyanNeon}{\thesubsection}}{0.7em}{}
  [\vspace{1pt}\color{grisLinea}\rule{\linewidth}{0.6pt}]
\titlespacing*{\subsection}{0pt}{14pt}{6pt}

\titleformat{\subsubsection}
  {\normalfont\normalsize\bfseries\color{azulNoche}}
  {\textcolor{cyanNeon}{\thesubsubsection}}{0.7em}{}
\titlespacing*{\subsubsection}{0pt}{10pt}{4pt}

% ── Cabeceras y pies de página ────────────────────────────────────────────────
\usepackage{fancyhdr}
\usepackage{lastpage}
\pagestyle{fancy}
\fancyhf{}
\renewcommand{\headrulewidth}{0pt}
\renewcommand{\footrulewidth}{0pt}
\fancyhead[L]{\textcolor{grisTexto}{\footnotesize\faIcon{robot}~Ecosistema Agente IA --- Destilación y Resolución Neuro-Simbólica}}
\fancyhead[R]{\textcolor{grisTexto}{\footnotesize\faIcon{calendar-alt}~Septiembre 2026}}
\fancyfoot[C]{%
  \begin{minipage}{\linewidth}
    \vspace{2pt}
    \color{cyanNeon}\rule{\linewidth}{1.2pt}\\[2pt]
    \centering\color{grisTexto}\footnotesize
    Página \thepage\ de \pageref{LastPage}
  \end{minipage}%
}

% ── Hiperenlaces ──────────────────────────────────────────────────────────────
\usepackage[colorlinks=true, linkcolor=azulNoche, urlcolor=cyanNeon, citecolor=azulMedio]{hyperref}
\sloppy
\emergencystretch 3em

% ── Paleta de colores — tecnológico dark-mode ──────────────────────────────────
\definecolor{fondoOscuro}{HTML}{0D1B2A}     % Dark navy — fondo banda título
\definecolor{verdeTurquesa}{HTML}{004D40}    % Verde turquesa — bandas de marca
\definecolor{azulNoche}{HTML}{1B3A6B}       % Midnight blue — secciones, marcos
\definecolor{azulMedio}{HTML}{2A5298}       % Azul medio — variante de acento
\definecolor{cyanNeon}{HTML}{00C8E0}        % Cyan eléctrico — acentos primarios
\definecolor{verdeSignal}{HTML}{00B887}     % Verde señal — fortalezas, éxito
\definecolor{naranjaVivo}{HTML}{FF7043}     % Naranja vivo — mejoras, advertencias
\definecolor{rojoAlerta}{HTML}{E53935}      % Rojo alerta — errores
\definecolor{grisPapel}{HTML}{F4F6FA}       % Fondo tarjetas claras
\definecolor{grisLinea}{HTML}{C8D0E0}       % Bordes sutiles
\definecolor{grisTexto}{HTML}{6B7A99}       % Texto secundario
\definecolor{amarilloNota}{HTML}{FFB300}    % Notas / advertencias doradas

% Colores específicos para bloques de código (dark-mode)
\definecolor{codigoFondo}{HTML}{1E2D3D}      % Fondo del bloque de código
\definecolor{codigoComentario}{HTML}{7A8FA6} % Comentarios
\definecolor{codigoCadena}{HTML}{FF9D5C}     % Cadenas / strings
\definecolor{codigoKeyword}{HTML}{00C8E0}    % Palabras clave
\definecolor{codigoNumero}{HTML}{00E5A0}     % Números

% ── Banda de título: portada infográfica ──────────────────────────────────────
\newcommand{\iconoBanda}{brain}
\newcommand{\bandaTitulo}[3][fondoOscuro]{%
  \begin{tcolorbox}[
    enhanced,
    colback=#1, colframe=#1,
    boxrule=0pt, arc=8pt, outer arc=8pt,
    width=\linewidth,
    top=16pt, bottom=14pt, left=14pt, right=14pt,
    borderline south={3.5pt}{0pt}{cyanNeon},
  ]
    \begin{minipage}[c]{0.08\linewidth}
      \centering
      \textcolor{cyanNeon}{\fontsize{32}{32}\selectfont\faIcon{\iconoBanda}}
    \end{minipage}%
    \hspace{10pt}%
    \begin{minipage}[c]{0.88\linewidth}
      {\fontsize{19}{23}\selectfont\bfseries\color{white}#2}\par\vspace{5pt}%
      \textcolor{cyanNeon!90!white}{\small\faIcon{angle-right}~#3}%
    \end{minipage}
    \vspace{2pt}
  \end{tcolorbox}%
  \vspace{6pt}%
}

% ── Tarjetas de datos (infografía) ────────────────────────────────────────────
\tcbset{
  tarjetaDato/.style={
    enhanced, breakable,
    arc=6pt, outer arc=6pt,
    colback=grisPapel, colframe=azulNoche,
    boxrule=1pt,
    fonttitle=\bfseries\small, coltitle=white,
    attach boxed title to top left={yshift=-3mm, xshift=6mm},
    boxed title style={
      colback=azulNoche, colframe=azulNoche,
      arc=4pt, boxrule=0pt,
    },
    drop fuzzy shadow=azulNoche!30!white,
    top=8pt, bottom=8pt, left=10pt, right=10pt,
  },
  tarjetaIcono/.style={
    enhanced, breakable,
    arc=6pt, outer arc=6pt,
    colback=white, colframe=grisLinea, boxrule=0.8pt,
    fonttitle=\bfseries\small, coltitle=azulNoche,
    borderline west={3pt}{0pt}{cyanNeon},
    drop fuzzy shadow=azulNoche!25!white,
    top=6pt, bottom=6pt, left=10pt, right=10pt,
  },
  cajaTitulo/.style={
    enhanced, breakable,
    arc=6pt, outer arc=6pt,
    colback=azulNoche!10!grisPapel, colframe=azulNoche,
    boxrule=1pt,
    fonttitle=\bfseries, coltitle=white,
    attach boxed title to top left={yshift=-3mm, xshift=6mm},
    boxed title style={
      colback=azulNoche, colframe=azulNoche,
      arc=4pt, boxrule=0pt,
    },
    drop fuzzy shadow=azulNoche!25!white,
    top=8pt, bottom=8pt, left=10pt, right=10pt,
  },
  cajaContenido/.style={
    enhanced, breakable,
    arc=5pt, outer arc=5pt,
    colback=grisPapel, colframe=grisLinea,
    boxrule=0.8pt,
    fonttitle=\bfseries\small, coltitle=azulNoche,
    top=6pt, bottom=6pt, left=10pt, right=10pt,
  },
  cajaFortaleza/.style={
    enhanced, breakable,
    arc=6pt, outer arc=6pt,
    colback=verdeSignal!8!white, colframe=verdeSignal,
    boxrule=1pt,
    borderline west={4pt}{0pt}{verdeSignal},
    fonttitle=\bfseries\small, coltitle=verdeSignal!70!black,
    drop fuzzy shadow=verdeSignal!20!white,
    top=6pt, bottom=6pt, left=10pt, right=10pt,
  },
  cajaMejora/.style={
    enhanced, breakable,
    arc=6pt, outer arc=6pt,
    colback=naranjaVivo!8!white, colframe=naranjaVivo,
    boxrule=1pt,
    borderline west={4pt}{0pt}{naranjaVivo},
    fonttitle=\bfseries\small, coltitle=naranjaVivo!80!black,
    drop fuzzy shadow=naranjaVivo!20!white,
    top=6pt, bottom=6pt, left=10pt, right=10pt,
  },
  cajaRecomendacion/.style={
    enhanced, breakable,
    arc=6pt, outer arc=6pt,
    colback=cyanNeon!6!white, colframe=cyanNeon!60!azulNoche,
    boxrule=1pt,
    borderline west={4pt}{0pt}{cyanNeon},
    fonttitle=\bfseries\small, coltitle=azulNoche,
    drop fuzzy shadow=azulNoche!20!white,
    top=6pt, bottom=6pt, left=10pt, right=10pt,
  },
}

% ── Ícono + texto ────────────────────────────────────────────────────────────
\newcommand{\iconotexto}[2]{\textcolor{cyanNeon}{\faIcon{#1}}~#2}

% ── Listados de código: estilo dark + auto-envoltura ─────────────────────────
\lstdefinestyle{estiloCodigo}{
    backgroundcolor=\color{codigoFondo},
    basicstyle=\ttfamily\footnotesize\color{white},
    commentstyle=\color{codigoComentario}\itshape,
    stringstyle=\color{codigoCadena},
    keywordstyle=\color{codigoKeyword}\bfseries,
    numberstyle=\tiny\color{codigoComentario},
    numbers=left,
    numbersep=8pt,
    showstringspaces=false,
    breaklines=true,
    breakatwhitespace=false,
    tabsize=4,
    frame=none,
    captionpos=b,
    aboveskip=4pt,
    belowskip=4pt,
    xleftmargin=12pt,
    framexleftmargin=12pt,
    literate=
      {á}{{\'a}}1 {é}{{\'e}}1 {í}{{\'i}}1 {ó}{{\'o}}1 {ú}{{\'u}}1
      {Á}{{\'A}}1 {É}{{\'E}}1 {Í}{{\'I}}1 {Ó}{{\'O}}1 {Ú}{{\'U}}1
      {ñ}{{\~n}}1 {Ñ}{{\~N}}1 {ü}{{\"u}}1 {Ü}{{\"U}}1
      {¿}{{?`}}1 {¡}{{!`}}1
}
\lstdefinestyle{estiloPython}{
    style=estiloCodigo,
    language=Python,
}
\lstdefinestyle{estiloBash}{
    style=estiloCodigo,
    language=bash,
}
\lstset{style=estiloCodigo}
\BeforeBeginEnvironment{lstlisting}{%
  \begin{tcolorbox}[
    enhanced, breakable,
    arc=6pt, outer arc=6pt,
    colback=codigoFondo, colframe=azulNoche,
    boxrule=1pt,
    drop fuzzy shadow=azulNoche!25!white,
    top=2pt, bottom=2pt, left=0pt, right=0pt,
  ]%
}
\AfterEndEnvironment{lstlisting}{\end{tcolorbox}}

% ── Cajas de contenido temáticas ──────────────────────────────────────────────
\newtcolorbox{cajaRecuerda}{
    enhanced, breakable,
    arc=6pt, outer arc=6pt,
    colback=azulNoche!10!grisPapel,
    colframe=azulNoche, boxrule=1pt,
    borderline west={4pt}{0pt}{cyanNeon},
    fonttitle=\bfseries\small\color{white},
    title={\faIcon{info-circle}~Recuerda},
    attach boxed title to top left={yshift=-3mm, xshift=6mm},
    boxed title style={colback=azulNoche, arc=4pt, boxrule=0pt},
    drop fuzzy shadow=azulNoche!25!white,
    left=10pt, right=10pt, top=8pt, bottom=8pt,
}

\newtcolorbox{cajaConcepto}{
    enhanced, breakable,
    arc=6pt, outer arc=6pt,
    colback=verdeSignal!8!white,
    colframe=verdeSignal, boxrule=1pt,
    borderline west={4pt}{0pt}{verdeSignal},
    fonttitle=\bfseries\small\color{white},
    title={\faIcon{lightbulb}~Concepto clave},
    attach boxed title to top left={yshift=-3mm, xshift=6mm},
    boxed title style={colback=verdeSignal!80!black, arc=4pt, boxrule=0pt},
    drop fuzzy shadow=verdeSignal!20!white,
    left=10pt, right=10pt, top=8pt, bottom=8pt,
}

\newtcolorbox{cajaEjemplo}{
    enhanced, breakable,
    arc=6pt, outer arc=6pt,
    colback=cyanNeon!5!white,
    colframe=cyanNeon!70!azulNoche, boxrule=1pt,
    borderline west={4pt}{0pt}{cyanNeon},
    fonttitle=\bfseries\small\color{fondoOscuro},
    title={\faIcon{code}~Ejemplo},
    attach boxed title to top left={yshift=-3mm, xshift=6mm},
    boxed title style={colback=cyanNeon, arc=4pt, boxrule=0pt},
    drop fuzzy shadow=azulNoche!20!white,
    left=10pt, right=10pt, top=8pt, bottom=8pt,
}

\newtcolorbox{cajaAlerta}{
    enhanced, breakable,
    arc=6pt, outer arc=6pt,
    colback=naranjaVivo!8!white,
    colframe=naranjaVivo, boxrule=1pt,
    borderline west={4pt}{0pt}{naranjaVivo},
    fonttitle=\bfseries\small\color{white},
    title={\faIcon{exclamation-triangle}~Punto crítico de diseño},
    attach boxed title to top left={yshift=-3mm, xshift=6mm},
    boxed title style={colback=naranjaVivo, arc=4pt, boxrule=0pt},
    drop fuzzy shadow=naranjaVivo!20!white,
    left=10pt, right=10pt, top=8pt, bottom=8pt,
}

% ── Documento Principal ───────────────────────────────────────────────────────
\begin{document}

% ── Portada Infográfica ───────────────────────────────────────────────────────
\bandaTitulo[fondoOscuro]{Destilación y Resolución con Skills Neuro-Simbólicos}{Fase 1: Síntesis de Libro a Skill \quad|\quad Fase 2: Razonamiento y Oráculo SymPy}

\begin{tcolorbox}[tarjetaDato, title={\faIcon{info-circle}~Ficha Técnica del Documento de Estudio}]
  \begin{tabularx}{\linewidth}{>{\bfseries\color{azulNoche}}l X >{\bfseries\color{azulNoche}}l X}
    \faIcon{calendar-check}~Fecha: & Septiembre 2026 & \faIcon{layer-group}~Arquitectura: & 3 Capas Deterministas \\
    \faIcon{coins}~Coste E2E (452 págs): & $\sim$\$0.75 USD (Fase 1 + Fase 2) & \faIcon{shield-alt}~Oráculo Simbólico: & SymPy 1.14.0 (Sandbox AST) \\
    \faIcon{microchip}~Modelos Tier: & DeepSeek V4 Pro / Gemini Flash & \faIcon{check-double}~Tasa Validación: & 33/33 Bloques SymPy (100\%) \\
  \end{tabularx}
\end{tcolorbox}

\vspace{-4pt}
\tableofcontents
\vspace{10pt}

% ═══════════════════════════════════════════════════════════════════════════════
% PARTE I: FASE 1 — SÍNTESIS DE LIBRO A SKILL
% ═══════════════════════════════════════════════════════════════════════════════
\newpage
\part*{PARTE I: Fase 1 --- De un Libro PDF a un Skill Estructurado}
\addcontentsline{toc}{part}{PARTE I: Fase 1 --- De un Libro PDF a un Skill Estructurado}

\section[Qué Resuelve la Fase 1]{\iconotexto{bullseye}{Qué Resuelve la Fase 1: Destilación Neuro-Simbólica}}

Un modelo de lenguaje (LLM) que simplemente ``responde preguntas sobre un libro técnico'' tiende a alucinar constantes, inventar ecuaciones y saltarse hipótesis fundamentales. La \textbf{Fase 1} transforma un documento técnico en PDF en un \textbf{skill estructurado} que permite al agente razonar formalmente:
\begin{itemize}[leftmargin=*]
  \item \textbf{Teoría y definiciones rigurosas:} Sin resúmenes superficiales; preservando supuestos y rango de validez.
  \item \textbf{Metodologías de resolución paso a paso:} Plantillas algorítmicas que replican el método del autor.
  \item \textbf{Límites de aplicabilidad y prerrequisitos:} Condiciones exactas de borde (\textit{cuándo NO aplicar}).
  \item \textbf{Fórmulas ejecutables en SymPy:} Cada fórmula matemática se representa como código ejecutable, permitiendo validación formal determinista con 0 alucinaciones.
\end{itemize}

\begin{tcolorbox}[cajaContenido, title={\faIcon{project-diagram}~Pipeline E2E de Destilación de la Fase 1}]
\begin{lstlisting}
[Libro PDF] --> extraer_libro_pdf.py --> texto_completo.txt --> enrutador.py
                                                                     |
                                             (destilacion chunk a chunk  |  -> corpus)
                                             (ensamblaje 7 archivos)    |
                                                                     v
                                              SKILL.md + references/ (formulas,
                                              metodologias, limites, prerrequisitos,
                                              tablas, glosario)
                                                                     v
                                        [validar_skill_formulas.py -> gate instalacion]
                                        [generar_latex_skill.py -> reporte tecnico PDF]
                                        [instalar_skill.py -> ~/.config/opencode/skills]
\end{lstlisting}
\end{tcolorbox}

% ── Sección 2: Arquitectura 3 Capas ───────────────────────────────────────────
\section[Arquitectura de 3 Capas (Fase 1)]{\iconotexto{layer-group}{Arquitectura de 3 Capas de la Fase 1}}

Siguiendo el principio de diseño del proyecto, toda la funcionalidad se divide estrictamente en 3 capas desacopladas:

\begin{table}[h!]
\centering
\small
\begin{tabularx}{\linewidth}{>{\bfseries\color{azulNoche}}l >{\ttfamily}p{5.2cm} >{\raggedright\arraybackslash}X}
\toprule
\rowcolor{azulNoche!12}\textbf{Capa} & \textbf{Archivo / Artefacto} & \textbf{Responsabilidad y Rol} \\
\midrule
\textbf{Directiva} & directives/libro\_a\_skill.yaml & SOP determinista: 8 pasos secuenciales, 13 edge cases catalogados y política de reintentos. \\
\midrule
\textbf{Orquestación} & flujo\_libro\_a\_skill.py & Entrevista al usuario, orquestación de scripts, manejo de estados en \texttt{.tmp/} y gate de instalación. \\
\midrule
\textbf{Ejecución} & execution/sintetizar\_skill.py \newline execution/extraer\_libro\_pdf.py \newline execution/validar\_skill\_formulas.py \newline execution/generar\_latex\_skill.py \newline execution/instalar\_skill.py & Scripts unitarios deterministas con entrada/salida estricta. Cero lógica de negocio en prompts del chat. \\
\bottomrule
\end{tabularx}
\caption{Matriz de responsabilidades de la arquitectura de 3 capas en la Fase 1.}
\end{table}

% ── Sección 3: Orquestador flujo_libro_a_skill.py ─────────────────────────────
\section[El Orquestador flujo\_libro\_a\_skill.py]{\iconotexto{cogs}{El Orquestador \texttt{flujo\_libro\_a\_skill.py}}}

El orquestador ejecuta una secuencia rigurosa de 8 pasos:

\begin{enumerate}[leftmargin=*]
  \item \textbf{Paso 1 --- Entrevista interactiva / CLI:} Captura nombre (\textit{snake\_case}), tema, idioma y alcance. Persiste estado en \path{.tmp/entrevista_skill_<nombre>.json}.
  \item \textbf{Paso 2 --- Extracción determinista (0 créditos):} Ejecuta \texttt{extraer\_libro\_pdf.py}. En el test E2E real extrajo 452 páginas, 1.6M caracteres y $\sim$400K tokens.
  \item \textbf{Paso 3 --- Enrutamiento inteligente:} Invoca \texttt{execution/enrutador.py --task contexto\_masivo} para seleccionar deterministamente el tier \textbf{deepseek} (\texttt{deepseek/deepseek-v4-pro}).
  \item \textbf{Paso 4 --- Síntesis y destilación (créditos):} Invoca \texttt{sintetizar\_skill.py} produciendo los 7 archivos estructurados.
  \item \textbf{Paso 5 --- Validación neuro-simbólica (0 créditos):} Ejecuta \texttt{validar\_skill\_formulas.py}. Si hay fórmulas inválidas, dispara bucle de reflexión hasta $\le 3$ veces.
  \item \textbf{Paso 6 --- Reporte LaTeX infográfico (0 créditos):} Genera reporte PDF en \path{docs/SKILL/<nombre>/}.
  \item \textbf{Paso 7 --- Instalación global condicionada (0 créditos):} Copia a \path{~/.config/opencode/skills/<nombre>/}.
  \item \textbf{Paso 8 --- Alerta sonora y telemetría:} Notifica al usuario con \texttt{alert\_user.py} avisando de reiniciar el entorno.
\end{enumerate}

\begin{cajaAlerta}
\textbf{Gate de Instalación Seguro (Implementado 2026-09-01):} Si la validación del Paso 5 arroja bloques simbólicos rotos, la instalación en el directorio global se \textbf{omite automáticamente} emitiendo un warning claro, pero \textbf{no se aborta el flujo}. De esta forma se genera siempre el reporte LaTeX con los diagnósticos sin contaminar el entorno de producción con skills defectuosos.
\end{cajaAlerta}

\subsection{Estrategia Determinista de Chunking}
La función \texttt{\_repartir\_chunks(texto, max\_tokens)} divide el texto por párrafos (\texttt{\textbackslash n\textbackslash n}) hasta alcanzar \texttt{max\_chars = max\_tokens * 4}. Párrafos más extensos que el límite se dividen a nivel de página. Con el default de \texttt{16000} tokens por chunk, un libro de 400K tokens se particiona limpiamente en $\approx 27$ chunks.

% ── Sección 4: Motor sintetizar_skill.py ──────────────────────────────────────
\section[Motor de Síntesis sintetizar\_skill.py]{\iconotexto{microchip}{Motor de Síntesis: \texttt{sintetizar\_skill.py}}}

\subsection{Dos Fases Internas de Procesamiento}
\begin{itemize}[leftmargin=*]
     \item \textbf{Fase A --- Destilación Chunk a Chunk (DeepSeek):} Cada chunk se envía con un prompt especializado que produce un resumen denso en markdown con definiciones, variables y ecuaciones en notación \texttt{\$...\$}. Fragmentos no técnicos devuelven \texttt{[sin contenido relevante]} y se filtran. El resultado se guarda en \path{.tmp/skill_<nombre>_destilado.txt} como \textbf{corpus destilado}.
  \item \textbf{Fase B --- Ensamblaje Estructurado (DeepSeek):} Genera un JSON único con \texttt{SKILL.md} y exactamente las 6 referencias formales obligatorias:
  \begin{enumerate}
    \item \texttt{formulas.md}: Ecuaciones con SymPy, dominio y notas de validez.
    \item \texttt{metodologias.md}: Plantillas numeradas paso a paso del procedimiento de cálculo.
    \item \texttt{limites\_aplicabilidad.md}: Supuestos de borde (\textit{cuándo NO aplicar}).
    \item \texttt{prerrequisitos.md}: Grafo de dependencias de conceptos y herramientas.
    \item \texttt{tablas.md} y \texttt{glosario.md}: Parámetros de componentes y definiciones formales.
  \end{enumerate}
\end{itemize}

\subsection{Saneo Determinista del JSON y Frontmatter}
Debido a imperfecciones típicas de LLMs al emitir JSONs de gran tamaño, se aplican algoritmos deterministas de reparación:
\begin{itemize}[leftmargin=*]
  \item \textbf{\texttt{\_sanear\_json}:} Corrige \textit{trailing commas} (\texttt{, \}} $\to$ \texttt{\}}) y escapa barras invertidas de LaTeX crudas (\texttt{\textbackslash Omega} $\to$ \texttt{\textbackslash\textbackslash Omega}) sin alterar caracteres de escape válidos.
  \item \textbf{\texttt{\_find\_balanced\_json}:} Extractor determinista mediante escaneo de llaves balanceadas que respeta cadenas y escapes, evitando fallos de expresiones regulares no codiciosas.
  \item \textbf{\texttt{\_normalizar\_frontmatter} + \texttt{\_yaml\_q}:} Re-emite escalares YAML entre comillas dobles, evitando que caracteres especiales como dos puntos (\texttt{:}) rompan el encabezado del skill.
\end{itemize}

\begin{cajaConcepto}
\textbf{Presupuesto de Ensamblaje y Solución de Truncamiento:} En libros completos, las 6 referencias estructuradas superan el límite de salida estándar de 8192 tokens. Se implementó el parámetro \texttt{--estructura-max-tokens 32768} para tier DeepSeek y el flag \texttt{--usar-corpus}, el cual permite re-ensamblar un skill existente sin volver a pagar la destilación chunk a chunk (ahorro de $\sim$30 min y $\sim$\$0.70 USD).
\end{cajaConcepto}

% ── Sección 5: Validación Neuro-Simbólica ─────────────────────────────────────
\section[Validación Neuro-Simbólica y Oráculo SymPy]{\iconotexto{shield-alt}{Validación Neuro-Simbólica: \texttt{validar\_skill\_formulas.py}}}

El validador escanea todos los bloques Python en \texttt{SKILL.md} y \texttt{references/*.md} sin consumir créditos de API:

\begin{enumerate}[leftmargin=*]
  \item \textbf{Análisis de Sintaxis:} Parseo de Abstract Syntax Tree (\texttt{ast.parse}).
  \item \textbf{Escaneo de Seguridad AST:} Bloquea llamadas inseguras (\texttt{os}, \texttt{subprocess}, \texttt{eval}, \texttt{exec}, dunders).
  \item \textbf{Ejecución en Sandbox Aislado:} Ejecución controlada con \texttt{subprocess.run([sys.executable, -c, src], timeout=5\,s)} en el validador de la Fase 1, y \texttt{timeout=30\,s} en el oráculo de la Fase 2.
  \item \textbf{Verificación de Estructuras Obligatorias:} Comprueba presencia de metodologías, límites y prerrequisitos.
\end{enumerate}

% ── Sección 6: Resultados Experimentales ──────────────────────────────────────
\section[Resultados Experimentales E2E]{\iconotexto{chart-bar}{Resultados Experimentales en Libro Completo (Fase 1)}}

Datos obtenidos en la prueba de validación E2E sobre el libro \textit{``Circuitos y Dispositivos Electrónicos''} (Lluís Prat Viñas):

\begin{table}[h!]
\centering
\small
\begin{tabularx}{\linewidth}{>{\bfseries\color{azulNoche}}l >{\raggedright\arraybackslash}X}
\toprule
\rowcolor{azulNoche!12}\textbf{Métrica de Evaluación} & \textbf{Valor Observado en Ejecución Real} \\
\midrule
Volumen de entrada & 452 páginas / 1.6M caracteres / $\sim$400K tokens \\
Modelo y Tier de enrutamiento & Tier \textbf{deepseek} (\texttt{deepseek/deepseek-v4-pro}) \\
Bloques SymPy extraídos y validados & \textbf{33/33 Bloques OK (100\% de éxito, 0 errores)} \\
Re-síntesis por reflexión requerida & 0 iteraciones adicionales (éxito en pasada inicial) \\
Artefactos estructurados producidos & 7 archivos (\texttt{SKILL.md} + 6 referencias técnicas) \\
Ubicación de reportes & \path{docs/SKILL/circuitos_dispositivos_electronicos/} \\
Instalación final & \path{~/.config/opencode/skills/circuitos_dispositivos_electronicos} \\
Coste total de la fase & $\sim$\$0.75 USD \\
\bottomrule
\end{tabularx}
\caption{Resultados del proceso de destilación neuro-simbólica E2E.}
\end{table}

% ═══════════════════════════════════════════════════════════════════════════════
% PARTE II: FASE 2 — RESOLUCIÓN DE PROBLEMAS
% ═══════════════════════════════════════════════════════════════════════════════
\newpage
\part*{PARTE II: Fase 2 --- Resolver Problemas con un Skill Estructurado}
\addcontentsline{toc}{part}{PARTE II: Fase 2 --- Resolver Problemas con un Skill Estructurado}

\section[Qué Resuelve la Fase 2]{\iconotexto{calculator}{Qué Resuelve la Fase 2: Patrón Neuro-Simbólico}}

Habiendo creado los skills en la Fase 1, la \textbf{Fase 2} los utiliza para resolver problemas cuantitativos de ingeniería. Soporta \textbf{múltiples skills del mismo dominio} (Opción C): el resolver hace retrieval sobre \emph{todos} ellos y combina los resultados, ampliando el campo de conocimiento sin re-destilar. El principio fundamental es:
\begin{center}
  \textit{El LLM jamás calcula mentalmente: razona cualitativamente y formula un programa simbólico; el oráculo determinista ejecuta y valida.}
\end{center}

\begin{tcolorbox}[cajaContenido, title={\faIcon{sync-alt}~Flujo de Razonamiento, Formulación y Validación Simbólica}]
\begin{lstlisting}
[Problema + N Skills del dominio]
        |
        v
[Retrieval Embeddings (multi-skill, 0 creditos)]
   - retrieval sobre TODOS los skills
   - combina y corta al top-k global
   - confianza_retrieval + mejor score por skill
        |
        v
[LLM: Formulador] --> [Oraculo SymPy]
   | (analisis +        |
   |  codigo SymPy)     |
   +---- feedback ------+
                |
      [Reflexion Local <=3]
\end{lstlisting}
\end{tcolorbox}

% ── Sección 8: Arquitectura 3 Capas (Fase 2) ──────────────────────────────────
\section[Arquitectura de 3 Capas (Fase 2)]{\iconotexto{layer-group}{Arquitectura de 3 Capas de la Fase 2}}

\begin{table}[h!]
\centering
\small
\begin{tabularx}{\linewidth}{>{\bfseries\color{azulNoche}}l >{\ttfamily}p{5.2cm} >{\raggedright\arraybackslash}X}
\toprule
\rowcolor{azulNoche!12}\textbf{Capa} & \textbf{Archivo / Artefacto} & \textbf{Responsabilidad y Rol} \\
\midrule
\textbf{Directiva} & directives/resolver\_skill.yaml & SOP: validación de entradas, pasos, 8 edge cases y presupuesto de retry. \\
\midrule
\textbf{Orquestación} & flujo\_resolver\_skill.py & Valida el/los problema(s) y skills, consulta al enrutador determinista, invoca el solver y reporta resultados. \\
\midrule
\textbf{Ejecución} & execution/resolver\_skill.py & Motor de retrieval multi-skill por similitud coseno, formulador LLM, sandbox del oráculo y reflexión. \\
\bottomrule
\end{tabularx}
\caption{Arquitectura de 3 capas para la resolución neuro-simbólica de problemas.}
\end{table}

% ── Sección 9: Motor resolver_skill.py ────────────────────────────────────────
\section[Motor de Ejecución resolver\_skill.py]{\iconotexto{cogs}{Motor de Ejecución: \texttt{resolver\_skill.py}}}

\subsection{1. Retrieval Semántico Multi-Skill por Embeddings (0 créditos)}
\begin{itemize}[leftmargin=*]
  \item Divide \texttt{SKILL.md} y referencias en secciones respetando encabezados markdown (\texttt{\# ... \#\#\#}).
  \item \textbf{Multi-skill (Opción C):} \texttt{--skill} acepta varios directorios de skills afines. El retrieval se ejecuta sobre \emph{todos} (cada uno con su caché de embeddings) y los resultados se \textbf{combinan con cuotas por skill}:
  \begin{itemize}[leftmargin=*]
    \item \textbf{Reserva obligatoria:} se incluye siempre la mejor sección de cada skill, para que un skill de dominio con pocas secciones de score alto nunca quede desplazado por un skill lateral con muchos chunks de score medio (evita dilución).
    \item \textbf{Cuotas equitativas:} el resto del $\text{top-}k$ se completa repartiendo lo más parejo posible entre los skills (round-robin por score desc), sin que un skill lateral sature el contexto.
    \item \textbf{Deduplicación:} la misma sección presente en varios skills se conserva una sola vez (la de mayor score).
    \item Cada sección se etiqueta con su \texttt{fuente} (skill de origen) y se reporta el \textbf{mejor score por skill} (\texttt{confianza\_retrieval.por\_skill}, calculado sobre todas sus secciones recuperadas, no solo las entregadas al LLM).
  \end{itemize}
  \item Genera representaciones vectoriales mediante el modelo local \texttt{paraphrase-multilingual-MiniLM-L12-v2}.
  \item Aplica similitud coseno seleccionando el $\text{top-}k$ de secciones con $\text{score} > \text{min\_score}$ (default $k=6$, umbral $0.05$).
  \item \textbf{Confianza del fundamento} (\texttt{confianza\_retrieval}): clasifica el \textbf{mejor score combinado} en \texttt{alta} ($\ge 0.55$), \texttt{media} ($\ge \text{alerta\_score}$, default $0.35$) o \texttt{baja} ($< 0.35$). Umbrales calibrados con evidencia medida (ver §10.2). Reporta además \texttt{mejor\_skill} (el más afín) y \texttt{skills\_sin\_secciones}.
  \item Implementa caché determinista en \path{.tmp/resolver_skill/<skill>/} con \texttt{secciones.json} y \texttt{embeddings.npz} basados en hashes \textbf{SHA-256} del contenido.
  \item \textbf{Ampliación del conocimiento:} si un problema es difícil y la confianza queda \texttt{baja}, ese es el indicador de que hace falta añadir más PDFs/skills del dominio (cada nuevo libro de la Fase 1 se usa como skill adicional sin re-destilar).
\end{itemize}

\subsection{2. Formulador LLM y Estructura de Salida}
El modelo recibe las secciones recuperadas y está obligado a responder con un JSON estricto de 3 claves:
\begin{tcolorbox}[cajaContenido, title={\faIcon{code}~Contrato de Datos del Formulador}]
\begin{lstlisting}
{
  "analisis": "Razonamiento paso a paso citando la Metodologia y limites",
  "codigo_sympy": "from sympy import *\n# Variables, ecuaciones y print formal",
  "resultado_esperado": "Valor numerico o simbolico esperado (prediccion)"
}
\end{lstlisting}
\end{tcolorbox}

\subsection{3. Oráculo SymPy y Bucle de Reflexión}
\begin{itemize}[leftmargin=*]
  \item \textbf{Seguridad previa:} Inspección AST con \texttt{\_escaneo\_seguridad}.
  \item \textbf{Sandbox:} Ejecución aislada con límite de tiempo (\texttt{timeout=30s}).
  \item \textbf{Bucle de Reflexión ($\le 3$):} Ante fallos de ejecución (\texttt{NameError}, \texttt{TypeError}, etc.), se inyecta el \texttt{stderr} real del compilador al LLM para corregir el código sin alterar la estrategia física de resolución.
\end{itemize}

% ── Sección 10: Validación Experimental Fase 2 ────────────────────────────────
\section[Validación Experimental Fase 2]{\iconotexto{flask}{Validación Experimental en Circuito BJT Real}}

Se evaluó la resolución del punto de operación Q de un transistor BJT en configuración emisor común con divisor de tensión en base ($V_{CC} = \SI{12}{\volt}$, $R_1 = \SI{100}{\kilo\ohm}$, $R_2 = \SI{50}{\kilo\ohm}$, $R_E = \SI{1}{\kilo\ohm}$, $R_C = \SI{2}{\kilo\ohm}$, $\beta = 100$):

\begin{tcolorbox}[tarjetaDato, title={\faIcon{check-circle}~Resultados de la Ejecución Experimental de Fase 2}]
  \begin{tabularx}{\linewidth}{>{\bfseries\color{azulNoche}}l X}
    Sección top recuperada: & \textit{``Metodología 7: Análisis BJT en DC''} (Score: \textbf{0.680}) \\
    Razonamiento del LLM: & Citó explícitamente la metodología del skill y verificó hipótesis activas de unión. \\
    Salida del Oráculo SymPy: & $I_{CQ} = \SI{2.4566}{\milli\ampere}, \quad V_{CEQ} = \SI{4.6057}{\volt}$ (Exactitud 100\%) \\
    Reflexiones necesarias: & \textbf{0 rondas} (Resuelto y validado a la primera pasada) \\
    Coste de inferencia: & $\sim$\$0.005 USD (utilizando Tier Flash para validación rápida) \\
  \end{tabularx}
\end{tcolorbox}

\subsection{Validación Experimental de Confianza del Retrieval (2026-09-03)}
Para que el sistema distinga robustamente \textit{problemas dentro de alcance} de los que
\textbf{no están reseñados en el PDF}, se calibaron los umbrales de confianza con evidencia
real medida sobre el skill del libro (modelo \texttt{paraphrase-multilingual-MiniLM-L12-v2}):

\begin{table}[h!]
\centering
\small
\begin{tabularx}{\linewidth}{>{\bfseries\color{azulNoche}}l >{\raggedright\arraybackslash}X}
\toprule
\rowcolor{azulNoche!12}\textbf{Tipo de problema} & \textbf{Mejor score coseno medido} \\
\midrule
Dentro de alcance (BJT punto Q, Ley de Ohm, Thévenin, Kirchhoff) & \textbf{0.45 -- 0.75} \\
Fuera de alcance (química orgánica, poesía griega) & \textbf{$\sim$0.25} \\
\bottomrule
\end{tabularx}
\caption{Separación empírica en/out-of-scope que sustenta la calibración de umbrales.}
\end{table}

Con los defaults iniciales (\texttt{alerta\_score}=0.20) la confianza \texttt{baja} jamás
se disparaba en la práctica (incluso un problema ajeno daba $\sim$0.25). Por ello se
**recalibró** con la separación medida: umbral de confianza débil \texttt{alerta\_score =
0.35} y tope de confianza \texttt{alta} en \texttt{0.55}.

\begin{tcolorbox}[tarjetaDato, title={\faIcon{shield-alt}~Detección de Fundamento Débil y Ampliación Multi-Skill (Medición Real, 0 créditos)}]
  \begin{tabularx}{\linewidth}{>{\bfseries\color{azulNoche}}l X}
    Problema fuera de alcance (química) & best score \textbf{0.2648} $\to$ confianza \textbf{\texttt{baja}} \\
    Problema dentro de alcance (BJT) & best score \textbf{0.7092} $\to$ confianza \textbf{\texttt{alta}} \\
    Abortar por fundamento débil (\texttt{--abortar-debil}) & exit \textbf{3} antes del LLM (0 créditos) \\
    Aviso al usuario & ``Confianza del retrieval BAJA (mejor score X)'' \\
    Formulador con fundamento débil & inyecta ``no cubierto por el libro'' en vez de forzar fórmula \\
    Multi-skill (Electrónica + Física) & mejor score por skill (p. ej. \texttt{0.68} / \texttt{0.35}); \texttt{mejor\_skill} = el más afín \\
    Confianza baja en problema difícil & señal de que \textbf{conviene añadir más PDFs/skills del dominio} \\
  \end{tabularx}
\end{tcolorbox}
El campo \texttt{confianza\_retrieval} \texttt{\{nivel, mejor\_score, alerta, alerta\_score,
skills\_consultados, skills\_sin\_secciones, mejor\_skill, por\_skill\}}
se incluye en el JSON de salida, permitiendo auditar la fiabilidad de la respuesta y
**detener el flujo (exit 3) sin gastar créditos** cuando el problema cae fuera del alcance
del knowledge de los skills. En modo multi-skill la confianza se calcula sobre el \textbf{mejor
score combinado} de todos los skills consultados.


% ═══════════════════════════════════════════════════════════════════════════════
% PARTE III: GUÍA DE COMANDOS Y COMPARATIVA GLOBAL
% ═══════════════════════════════════════════════════════════════════════════════
\newpage
\part*{PARTE III: Operatoria, Costes y Síntesis Global}
\addcontentsline{toc}{part}{PARTE III: Operatoria, Costes y Síntesis Global}

\section[Guía de Comandos Operativos]{\iconotexto{terminal}{Guía de Comandos Operativos}}

\begin{tcolorbox}[cajaContenido, title={\faIcon{terminal}~Comandos para la Fase 1 (Creación del Skill)}]
\begin{lstlisting}[style=estiloBash]
# 1. Flujo completo de sintesis con validacion y reporte
python3 flujo_libro_a_skill.py --pdf "libro.pdf" --nombre mi_skill \
    --tema "Electronica Analogica" --idioma es --sobrescribir --reflexion 1

# 2. Modo Dry-Run (solo genera y valida, sin instalar globalmente)
python3 flujo_libro_a_skill.py --pdf "libro.pdf" --nombre mi_skill \
    --tema "Electronica Analogica" --dry-run

# 3. Re-ensamblar reutilizando el corpus ya destilado (Ahorro de ~90% creditos)
python3 flujo_libro_a_skill.py --pdf "libro.pdf" --nombre mi_skill \
    --tema "Electronica Analogica" --sobrescribir --usar-corpus --estructura-max-tokens 32768

# 4. Validar formulas simbolicas de forma 100% local (0 creditos)
python3 execution/validar_skill_formulas.py --skill ~/.config/opencode/skills/mi_skill
\end{lstlisting}
\end{tcolorbox}

\begin{tcolorbox}[cajaContenido, title={\faIcon{terminal}~Comandos para la Fase 2 (Resolución de Problemas)}]
\begin{lstlisting}[style=estiloBash]
# 1. Resolver problema en modo produccion (calculo_formal -> Opus)
python3 flujo_resolver_skill.py --problema "Calcular el punto Q de un BJT..." \
    --skill ~/.config/opencode/skills/circuitos_dispositivos_electronicos

# 2. Test rapido y economico con override a modelo Flash (~$0.005)
python3 flujo_resolver_skill.py --problema "Calcular el punto Q de un BJT..." \
    --skill ~/.config/opencode/skills/circuitos_dispositivos_electronicos \
    --modelo google/gemini-3.7-flash --max-reflexion 1

# 3. Inspeccionar unicamente las secciones relevantes recuperadas (0 creditos)
python3 execution/resolver_skill.py --problema "Calcular el punto Q de un BJT..." \
    --skill ~/.config/opencode/skills/circuitos_dispositivos_electronicos --solo-retrieval

# 4. Abortar si el retrieval tiene fundamento debil (confianza baja) -> exit 3
python3 flujo_resolver_skill.py --problema "..." \
    --skill ~/.config/opencode/skills/circuitos_dispositivos_electronicos --abortar-debil

# 5. Ajustar el umbral de confianza debil (default 0.35; mas estricto = mas alto)
python3 flujo_resolver_skill.py --problema "..." --skill <dir> --alerta-score 0.40

# 6. AMPLIAR el campo de conocimiento: varios skills del mismo dominio
#    (Electronica + Fisica). Retrieval combinado, top-k global, 0 creditos extra.
python3 flujo_resolver_skill.py --problema "Calcular el punto Q de un BJT..." \
    --skill ~/.config/opencode/skills/circuitos_dispositivos_electronicos \
    ~/.config/opencode/skills/electronica_2 ~/.config/opencode/skills/fisica_electronica

# 7. Solo retrieval multi-skill: ver que aporta cada skill (0 creditos)
python3 execution/resolver_skill.py --skill <dir1> <dir2> \
    --problema "..." --solo-retrieval
\end{lstlisting}
\end{tcolorbox}

\section[Comparativa y Síntesis Global]{\iconotexto{balance-scale}{Cuadro Comparativo: Fase 1 vs Fase 2}}

\begin{table}[h!]
\centering
\small
\begin{tabularx}{\linewidth}{>{\bfseries\color{azulNoche}}l >{\raggedright\arraybackslash}X >{\raggedright\arraybackslash}X}
\toprule
\rowcolor{azulNoche!12}\textbf{Eje de Comparación} & \textbf{Fase 1: Síntesis de Skill} & \textbf{Fase 2: Resolución con Skill} \\
\midrule
\textbf{Objetivo Primario} & Extraer y destilar el conocimiento de un libro en un skill formal. & Aplicar el conocimiento estructurado para resolver problemas. \\
\midrule
\textbf{Entrada} & Archivo PDF técnico (centenares de páginas). & Enunciado del problema de ingeniería. \\
\midrule
\textbf{Salida} & \texttt{SKILL.md} + 6 archivos en \texttt{references/}. & JSON con análisis, código SymPy y resultado verificado. \\
\midrule
\textbf{Rol del LLM} & Destilador y estructurador formal de contenido. & Razonador cualitativo y formulador simbólico. \\
\midrule
\textbf{Oráculo Determinista} & Validador estático/ejecutable de bloques SymPy. & Ejecutor dinámico del cálculo y evaluador de solución. \\
\midrule
\textbf{Consumo de Créditos} & Pago único por libro ($\sim$\$0.75 USD). & Fraccional por problema resuelto ($\sim$\$0.005 a \$0.03 USD). \\
\midrule
\textbf{Ahorro Inteligente} & Persistencia de corpus destilado (\texttt{--usar-corpus}). & Caché local de embeddings (\texttt{embeddings.npz}). \\
\bottomrule
\end{tabularx}
\caption{Matriz comparativa integral de las Fases 1 y 2 del ecosistema de IA.}
\end{table}

\begin{cajaRecuerda}
El enfoque neuro-simbólico desacopla la potencia de abstracción lingüística del LLM de la precisión matemática. Al forzar al modelo a emitir código SymPy verificado en sandbox, se eliminan por diseño las alucinaciones numéricas y se garantiza la reproducibilidad técnica exigida en ingeniería.
\end{cajaRecuerda}

\end{document}
